%% kapa: exam=1 round=28 year=2017 period=1 form=A subject=real_estate subject_ko=부동산학원론 seq=23 number=103
%% source: Q-Net 감정평가사 기출문제 (cst00309) artlSeq=5199603 file=extracted/1008/5199603/2017년 제28회 감정평가사 1차시험 문제지/1교시 A형.pdf
%% conversion: kordoc 4.24.51 -> tools/gen1.py (auto); review=proofread-sonnet; flags=none
\begin{question}{103}
감가수정에 관한 설명으로 옳은 것은?
\begin{choices}
  \item 치유가능한 감가는 내용년수 항목 중에서 치유로 증가가 예상되는 효용이 치유에 요하는 비용보다 큰 경우의 감가를 의미한다.
  \item 감가수정의 방법은 직접법과 간접법이 있으며, 직접법에는 내용년수법, 관찰감가법 및 분해법이 있다. 감가수정액의 산정은 이 세가지 방법을 병용하여 산정해야 한다.
  \item 감가수정은 재조달원가에서 부동산가격에 영향을 미치는 물리적ㆍ기능적ㆍ경제적 감가요인 등을 고려하고, 그에 해당하는 감가수정액을 공제하여, 기준시점 현재 대상물건의 기간손익의 배분을 산정하기 위한 것이다.
  \item 감정평가대상이 되는 부동산의 상태를 면밀히 관찰한 후 감정평가사의 폭넓은 경험과 지식에 의존하는 것이 분해법이다.
  \item 감가요인을 물리적ㆍ기능적ㆍ경제적 요인으로 세분하고, 치유가능ㆍ불능항목으로 세분하여 각각의 발생감가의 합계액을 감가수정액으로 하는 방법이 관찰감가법이다.
\end{choices}
\end{question}
\answer{1}
